\textbf{Phần 1. Trắc nghiệm lựa chọn}
\setcounterpageref{ex}{0}

\begin{ex} %[3X3V81]
Tính giá trị của biểu thức: $A = 25\,000 \times 3 - 48\,000 : 6$.
\choice
{$65\,000$}
{$70\,000$}
{\True $67\,000$}
{$66\,000$}
\loigiai{
$25\,000 \times 3 = 75\,000$; $48\,000 : 6 = 8000$; $75\,000 - 8000 = 67\,000$.
\textbf{Đáp án của câu hỏi là C}
}
\end{ex}

\begin{ex} %[3X3V81]
Số liền sau của số lớn nhất có năm chữ số là
\choice
{\True $100\,000$}
{$99\,998$}
{$100\,001$}
{$99\,990$}
\loigiai{
$99\,999 + 1 = 100\,000$.
\textbf{Đáp án của câu hỏi là A}
}
\end{ex}

\begin{ex} %[3X3V81]
Một hình chữ nhật có chiều dài $30\text{ cm}$, chiều rộng bằng $\dfrac{1}{3}$ chiều dài. Diện tích của hình chữ nhật đó là
\choice
{$100\text{ cm}^2$}
{$200\text{ cm}^2$}
{$400\text{ cm}^2$}
{\True $300\text{ cm}^2$}
\loigiai{
Chiều rộng: $30 : 3 = 10\text{ cm}$. Diện tích: $30 \times 10 = 300\text{ (cm}^2)$.
\textbf{Đáp án của câu hỏi là D}
}
\end{ex}

\begin{ex} %[3X3V81]
Một hình vuông có chu vi $64\text{ cm}$. Diện tích của hình vuông đó là
\choice
{$128\text{ cm}^2$}
{\True $256\text{ cm}^2$}
{$64\text{ cm}^2$}
{$512\text{ cm}^2$}
\loigiai{
Cạnh hình vuông: $64 : 4 = 16\text{ cm}$. Diện tích: $16 \times 16 = 256\text{ (cm}^2)$.
\textbf{Đáp án của câu hỏi là B}
}
\end{ex}

\begin{ex} %[3X3V81]
Số La Mã $\text{XXI}$ cộng với $\text{IV}$ được kết quả là số La Mã nào?
\choice
{$\text{XXIV}$}
{$\text{XXVI}$}
{\True $\text{XXV}$}
{$\text{XX}$}
\loigiai{
$\text{XXI} = 21$; $\text{IV} = 4$. $21 + 4 = 25 = \text{XXV}$.
\textbf{Đáp án của câu hỏi là C}
}
\end{ex}

\begin{ex} %[3X3V81]
Một cửa hàng có $80\,000\text{ kg}$ gạo. Ngày đầu bán $\dfrac{1}{4}$ số gạo, ngày thứ hai bán $\dfrac{1}{3}$ số gạo còn lại. Cửa hàng còn lại bao nhiêu ki-lô-gam gạo?
\choice
{\True $40\,000\text{ kg}$}
{$20\,000\text{ kg}$}
{$50\,000\text{ kg}$}
{$30\,000\text{ kg}$}
\loigiai{
Ngày 1 bán: $80\,000 : 4 = 20\,000\text{ kg}$. Còn lại: $80\,000 - 20\,000 = 60\,000\text{ kg}$. Ngày 2 bán: $60\,000 : 3 = 20\,000\text{ kg}$. Còn lại: $60\,000 - 20\,000 = 40\,000\text{ (kg)}$.
\textbf{Đáp án của câu hỏi là A}
}
\end{ex}

\begin{ex} %[3X3V81]
Làm tròn số $78\,650$ đến hàng nghìn ta được số
\choice
{$78\,000$}
{$80\,000$}
{$78\,700$}
{\True $79\,000$}
\loigiai{
Chữ số hàng trăm là 6 ($6 \ge 5$) nên làm tròn lên $79\,000$.
\textbf{Đáp án của câu hỏi là D}
}
\end{ex}

\begin{ex} %[3X3V81]
Một hình tròn có bán kính $9\text{ cm}$. Đoạn thẳng dài nhất nối 2 điểm trên hình tròn có độ dài là
\choice
{$9\text{ cm}$}
{\True $18\text{ cm}$}
{$27\text{ cm}$}
{$36\text{ cm}$}
\loigiai{
Đoạn dài nhất chính là đường kính: $9 \times 2 = 18\text{ (cm)}$.
\textbf{Đáp án của câu hỏi là B}
}
\end{ex}

\begin{ex} %[3X3V81]
Đổi $5\text{ m } 4\text{ cm}$ ra đơn vị mi-li-mét:
\choice
{$540\text{ mm}$}
{$5004\text{ mm}$}
{\True $5040\text{ mm}$}
{$5400\text{ mm}$}
\loigiai{
$5\text{ m} = 5000\text{ mm}$; $4\text{ cm} = 40\text{ mm}$. Tổng: $5040\text{ mm}$.
\textbf{Đáp án của câu hỏi là C}
}
\end{ex}

\begin{ex} %[3X3V81]
Tháng nào sau đây có đúng 30 ngày?
\choice
{\True Tháng 4}
{Tháng 1}
{Tháng 3}
{Tháng 5}
\loigiai{
Tháng 4 có đúng 30 ngày.
\textbf{Đáp án của câu hỏi là A}
}
\end{ex}

\begin{ex} %[3X3V81]
Một con xúc xắc gieo ngẫu nhiên. Sự kiện 'Xuất hiện mặt có số chấm bé hơn 7' là sự kiện
\choice
{Có thể xảy ra}
{Không thể xảy ra}
{Hiếm khi xảy ra}
{\True Chắc chắn xảy ra}
\loigiai{
Tất cả các mặt từ 1 đến 6 đều bé hơn 7 nên chắc chắn xảy ra.
\textbf{Đáp án của câu hỏi là D}
}
\end{ex}

\begin{ex} %[3X3V81]
Tổng của hai số là 90 000, số lớn gấp 4 lần số bé. Số lớn là
\choice
{$18\,000$}
{\True $72\,000$}
{$60\,000$}
{$75\,000$}
\loigiai{
Tổng số phần: $4 + 1 = 5$ phần. Số bé: $90\,000 : 5 = 18\,000$. Số lớn: $18\,000 \times 4 = 72\,000$.
\textbf{Đáp án của câu hỏi là B}
}
\end{ex}

\textbf{Phần 2. Trắc nghiệm đúng sai}
\setcounterpageref{ex}{0}

\begin{ex} %[3X3V81]
Xét tính đúng sai của các kiến thức tổng hợp toàn bộ chương trình Toán 3:
\choiceTF
{\True Số 100 000 là số tròn trăm nghìn bé nhất có sáu chữ số.}
{$1\text{ m} = 100\text{ mm}$.}
{\True Hình vuông có 4 cạnh bằng nhau và 4 góc vuông.}
{Trong phép chia cho 5, số dư lớn nhất có thể là 5.}
\loigiai{
\begin{itemize}
    \item Ý a) Đúng: Định nghĩa.
    \item Ý b) Sai: Sai vì $1\text{ m} = 1000\text{ mm}$.
    \item Ý c) Đúng: Đặc điểm hình vuông.
    \item Ý d) Sai: Sai vì số dư lớn nhất là 4.
\end{itemize}
\begin{itemize}
    \item a) \textbf{Đúng}.
    \item b) \textbf{Sai}.
    \item c) \textbf{Đúng}.
    \item d) \textbf{Sai}.
\end{itemize}
}
\end{ex}

\begin{ex} %[3X3V81]
Một hình chữ nhật có chiều dài $20\text{ cm}$, chiều rộng $10\text{ cm}$:
\choiceTF
{\True Chu vi hình chữ nhật là $60\text{ cm}$.}
{\True Diện tích hình chữ nhật là $200\text{ cm}^2$.}
{Nếu chia đôi hình chữ nhật theo chiều dài thì được 2 hình vuông cạnh $5\text{ cm}$.}
{\True Chiều dài gấp đôi chiều rộng.}
\loigiai{
\begin{itemize}
    \item Ý a) Đúng: $(20 + 10) \times 2 = 60\text{ cm}$.
    \item Ý b) Đúng: $20 \times 10 = 200\text{ cm}^2$.
    \item Ý c) Sai: Sai vì cắt đôi được 2 hình vuông cạnh $10\text{ cm}$.
    \item Ý d) Đúng: $20 : 10 = 2$ lần.
\end{itemize}
\begin{itemize}
    \item a) \textbf{Đúng}.
    \item b) \textbf{Đúng}.
    \item c) \textbf{Sai}.
    \item d) \textbf{Đúng}.
\end{itemize}
}
\end{ex}

\textbf{Phần 3. Trả lời ngắn}
\setcounterpageref{ex}{0}

\begin{ex} %[3X3V81]
Tính giá trị biểu thức: $15\,000 \times 4 + 40\,000 : 8$.
\shortans{65000}
\loigiai{
\textbf{Đáp số:} $65000$\\
$60\,000 + 5000 = 65\,000$.
}
\end{ex}

\begin{ex} %[3X3V81]
Tìm số liền sau của số 99 999.
\shortans{100000}
\loigiai{
\textbf{Đáp số:} $100000$\\
$99\,999 + 1 = 100\,000$.
}
\end{ex}

\begin{ex} %[3X3V81]
Một hình vuông có chu vi $72\text{ cm}$. Tính độ dài cạnh hình vuông theo xăng-ti-mét.
\shortans{18}
\loigiai{
\textbf{Đáp số:} $18$\\
$72 : 4 = 18\text{ (cm)}$.
}
\end{ex}

\begin{ex} %[3X3V81]
Tính diện tích hình chữ nhật có chiều dài $15\text{ cm}$ và chiều rộng $8\text{ cm}$.
\shortans{120}
\loigiai{
\textbf{Đáp số:} $120$\\
$15 \times 8 = 120\text{ (cm}^2)$.
}
\end{ex}

\begin{ex} %[3X3V81]
Đổi $2\ell - 750\text{ ml}$ ra mi-li-lít.
\shortans{1250}
\loigiai{
\textbf{Đáp số:} $1250$\\
$2000 - 750 = 1250\text{ (ml)}$.
}
\end{ex}

\begin{ex} %[3X3V81]
Tổng hai số là 80 000, hiệu hai số là 20 000. Tìm số lớn.
\shortans{50000}
\loigiai{
\textbf{Đáp số:} $50000$\\
Số lớn là: $(80\,000 + 20\,000) : 2 = 50\,000$.
}
\end{ex}
